\documentclass[11pt]{amsart}

% Encoding and fonts
\usepackage[T1]{fontenc}
\usepackage[utf8]{inputenc}
\usepackage{lmodern}

% Math packages
\usepackage{amsmath}
\usepackage{amssymb}
\usepackage{amsfonts}
\usepackage{amsthm}
\usepackage{mathtools}

% Layout and utilities
\usepackage{geometry}
\usepackage{enumitem}
\usepackage{booktabs}
\usepackage{array}
\usepackage{longtable}
\usepackage{microtype}
\geometry{margin=1in}

\numberwithin{equation}{section}
\setcounter{tocdepth}{3}
\setcounter{secnumdepth}{3}

% Theorem-like environments
\theoremstyle{plain}
\newtheorem{theorem}{Theorem}[section]
\newtheorem{proposition}[theorem]{Proposition}
\newtheorem{lemma}[theorem]{Lemma}
\newtheorem{corollary}[theorem]{Corollary}
\newtheorem{conjecture}[theorem]{Conjecture}

\theoremstyle{definition}
\newtheorem{definition}[theorem]{Definition}
\newtheorem{experiment}[theorem]{Experiment}
\newtheorem{heuristic}[theorem]{Heuristic}
\newtheorem{question}[theorem]{Question}
\newtheorem{example}[theorem]{Example}

\theoremstyle{remark}
\newtheorem{remark}[theorem]{Remark}
\newtheorem{warning}[theorem]{Warning}

% Notation
\newcommand{\Mert}{M}
\newcommand{\Mob}{\mu}
\newcommand{\Sq}{Q}
\newcommand{\eps}{\varepsilon}
\newcommand{\R}{\mathbb R}
\newcommand{\C}{\mathbb C}
\newcommand{\Ftwo}{\mathbb F_2}
\newcommand{\Cat}{\mathrm{Cat}}
\newcommand{\rank}{\operatorname{rank}}
\newcommand{\supp}{\operatorname{supp}}
\newcommand{\Cut}{\operatorname{Cut}}
\newcommand{\ord}{\operatorname{ord}}
\newcommand{\LIT}{\textup{[LIT]}}
\newcommand{\DED}{\textup{[DED]}}
\newcommand{\COMP}{\textup{[COMP]}}
\newcommand{\OPEN}{\textup{[OPEN]}}
\newcommand{\NEG}{\textup{[NEG]}}
\newcommand{\HEUR}{\textup{[HEUR]}}

\title[Arithmetic Selection and Catalan Geometry]{Arithmetic Selection and Catalan Geometry of the Mertens Walk:\\
Exact Structures, Capacity Deficits, Occupation Energy, and Open Routes}
\author{V\'ictor Casa\~na-Giner}
\address{Bio Logic Crop Science S.L., Spain}
\thanks{ORCID: 0000-0002-8252-2608}
\date{15 September 2026}

% hyperref must be loaded last
\usepackage[hidelinks]{hyperref}

\begin{document}
\begin{abstract}
We organize several exact and experimental structures around the M\"obius function
$\Mob(n)$ and its summatory function $\Mert(x)$.  The squarefree-compressed Mertens
walk is an exact nearest-neighbour path; complete zero-to-zero excursions are
primitive Dyck paths and open pieces are ballot meanders.  This produces a canonical
Catalan state space, but ambient Catalan rarity alone does not restrict the arithmetic
selector.  We therefore emphasize arithmetic mechanisms that can act on the whole
prefix profile.

The first part develops the occupation-energy formulation
$\mathcal E_q(x,H)=\sum_{h\le H}|\Mert(x+h)-\Mert(x)|^{2q}$ and the exact quadratic
operator $J_H^*J_H$.  The all-$q$ dyadic hierarchy is an exact RH criterion, while any
fixed power saving over the ballistic quadratic scale is a strictly weaker and
quantitatively meaningful first target.  The second part combines adversarial
multiplicative completions, signed displacement cuts, and collective Catalan capacities.
For a finite squarefree prime product $P$, the squarefree integers decompose into
components $aD(P)$, and the exact worst-case prefix capacity is expressed by an
$L^1$ norm of a signed divisor-window kernel.  Adding a prime acts by a delayed
anti-copy, gives a monotone capacity ratio, and also yields a lower barrier showing a
tradeoff between confinement and remaining entropy.  A further no-go theorem shows
that static Ramar\'e relations over a finite prime set do not couple different
$aD(P)$ components; genuine inter-component gain must use causal truncation or an
analytic recombination estimate.

The synthesis also retains the signed-cut extinction programme, the ordinal-prime
transform and its coercive negative-excursion energy, the divisibility-surface and
$p$-adic shell coordinates, and a one-way Catalan calibration of hypothetical
off-critical zeta zeros.  Finite computations through $N=455{,}052{,}511$ are recorded
only as experiments.  The resulting research frontier is explicit: obtain a proved,
non-tautological, sign-sensitive collective restriction that reduces high Catalan
sectors or occupation energy by a fixed power while leaving substantial arithmetic
freedom.
\end{abstract}

\maketitle
\tableofcontents
\bigskip

\begin{remark}[Status and purpose]
\label{rem:status-purpose}
This document is a research synthesis and working notebook.  It consolidates the strongest exact constructions, negative results, finite experiments, and open mechanisms developed in a collection of M\"obius--Mertens manuscripts and project notebooks.  The labels \LIT, \DED, \COMP, \OPEN, \NEG, and \HEUR mean, respectively: established literature input; rigorous deduction recorded in the project; finite computation; unresolved target; rigorous or structural negative result; and heuristic.  No literature-wide priority claim is made for deductions marked \DED without a separate bibliographic audit.  In particular, exact reformulations equivalent to the Riemann hypothesis are not presented as progress toward its proof unless an independent arithmetic estimate is supplied.
\end{remark}



\section{Main objects and logical boundary}
\label{sec:main-objects}

Let
\begin{equation}
\Mert(x)=\sum_{n\le x}\Mob(n),
\label{eq:mertens-definition}
\end{equation}
and let
\begin{equation}
\Sq(x)=\sum_{n\le x}\Mob(n)^2
=\#\{n\le x:n\text{ is squarefree}\}.
\label{eq:squarefree-count}
\end{equation}
The standard squarefree counting asymptotic is
\begin{equation}
\Sq(x)=\frac{6}{\pi^2}x+O(\sqrt x).
\label{eq:squarefree-asymptotic}
\end{equation}
The classical Mertens criterion states that
\begin{equation}
\mathrm{RH}
\quad\Longleftrightarrow\quad
\Mert(x)=O_{\eps}(x^{1/2+\eps})
\quad\text{for every }\eps>0;
\label{eq:mertens-rh-criterion}
\end{equation}
see, for example, \cite{Titchmarsh1986}.  The obsolete Mertens conjecture
$|\Mert(x)|<\sqrt x$ is false \cite{OdlyzkoTeRiele1985}; it must never be used as an
asymptotic surrogate for \eqref{eq:mertens-rh-criterion}.

\begin{warning}[Criterion versus mechanism]
\label{warn:criterion-mechanism}
Any construction that immediately reduces to \eqref{eq:mertens-rh-criterion}, to a
zero-free half-plane for $1/\zeta$, or to an equivalent mean-square statement is a
criterion or change of coordinates.  It becomes a proof route only after a genuinely
independent arithmetic estimate is established.
\end{warning}

The central methodological boundary used throughout this notebook is
\begin{equation}
\boxed{
\text{partial arithmetic information}
\Longrightarrow
\text{uniform non-tautological geometric or energetic inequality}.}
\label{eq:boundary-principle}
\end{equation}
The left side must not contain the desired Mertens bound in disguised form.

\section{The squarefree-compressed Mertens walk}
\label{sec:compressed-walk}

Let
\begin{equation}
1=q_1<q_2<q_3<\cdots
\label{eq:squarefree-order}
\end{equation}
be the squarefree integers in increasing order, and define
\begin{equation}
X_m=\sum_{j\le m}\Mob(q_j).
\label{eq:compressed-walk}
\end{equation}

\begin{proposition}[Exact compression \DED]
\label{prop:compression}
For every $m\ge1$ and real $x\ge1$,
\begin{align}
X_m&=\Mert(q_m),
\label{eq:compression-at-squarefree}\\
\Mert(x)&=X_{\Sq(x)}.
\label{eq:decompression}
\end{align}
\end{proposition}

\begin{proof}
The terms deleted from ordinary time are exactly the terms for which $\Mob(n)=0$.
Thus the nonzero increments are unchanged and only holding times are removed.
\end{proof}

Let
\begin{equation}
0=z_0<z_1<z_2<\cdots
\label{eq:return-times}
\end{equation}
be successive return times of $X_m$ to zero, when they exist.  Put
\begin{equation}
L_r=z_r-z_{r-1}=2n_r.
\label{eq:excursion-length}
\end{equation}
Parity forces each complete excursion length to be even.

\begin{proposition}[Primitive Catalan factorization \DED]
\label{prop:catalan-factorization}
Fix a finite return skeleton and the sign of each complete excursion.  The number of
nearest-neighbour paths realizing excursions of lengths $2n_1,\ldots,2n_R$ is
\begin{equation}
\mathcal N(z_1,\ldots,z_R)=\prod_{r=1}^{R}C_{n_r-1},
\label{eq:catalan-product}
\end{equation}
where
\begin{equation}
C_n=\frac{1}{n+1}\binom{2n}{n}
\label{eq:catalan-number}
\end{equation}
is the $n$th Catalan number.
\end{proposition}

\begin{proof}
After reflection of a negative excursion, removing the first up-step and final
down-step gives an arbitrary Dyck path of semilength $n_r-1$.  Distinct complete
excursions are combinatorially independent once the skeleton and signs are fixed.
\end{proof}

For open pieces the exact ballot count is equally important.

\begin{proposition}[Strict meander endpoint count \LIT]
\label{prop:meander-count}
Let $t\ge1$ and $h>0$ with $t\equiv h\pmod 2$.  The number of $\pm1$ paths from $0$
to $h$ in $t$ steps that remain strictly positive after time zero is
\begin{equation}
\#\mathcal M_{t,h}=\frac{h}{t}\binom{t}{(t+h)/2}.
\label{eq:meander-count}
\end{equation}
\end{proposition}

The ordinary-time embedding is retained by the holding times of the squarefree clock.
For a window $(a,b]$, let
\begin{equation}
w_j(a,b)=\#\{n\in(a,b]:\Sq(n)=j\}.
\label{eq:holding-weights}
\end{equation}
Then any ordinary-time occupation statistic can be written as a weighted statistic of
the compressed path.  Deleting zeros therefore preserves the path shape but not the
physical time at which a given height occurs.

\section{Catalan height as a calibrated geometric scale}
\label{sec:catalan-height}

Let $\mathcal D_n$ be the set of Dyck paths of semilength $n$ and let
$D_{n,h}$ denote the number whose height is at most $h$.

\begin{proposition}[Bounded-height spectral formula \LIT]
\label{prop:bounded-height}
For $n\ge1$ and $h\ge1$,
\begin{equation}
D_{n,h}
=
\frac{2}{h+2}
\sum_{j=1}^{h+1}
\sin^2\!\left(\frac{\pi j}{h+2}\right)
\left(2\cos\frac{\pi j}{h+2}\right)^{2n}.
\label{eq:bounded-height-spectral}
\end{equation}
\end{proposition}

This is the transfer-matrix formula for the path graph on heights
$0,1,\ldots,h$.  The corresponding Brownian-excursion scaling and moderate deviations
are classical; see \cite{Kaigh1976,DurrettIglehartMiller1977,KhorunzhiyMarckert2009,Kousha2012}.

\begin{proposition}[Height phase transition \LIT]
\label{prop:height-transition}
Fix $c>0$ and put $h_n=\lfloor c n^\beta\rfloor$.  Then the Catalan height threshold
occurs at $\beta=1/2$: for $\beta<1/2$,
\begin{equation}
\log\frac{D_{n,h_n}}{C_n}
=
-\frac{\pi^2}{c^2}n^{1-2\beta}+O(\log n),
\label{eq:subcritical-height}
\end{equation}
whereas for $\beta>1/2$ the ratio tends to one.  At $\beta=1/2$ the ratio has a
non-degenerate Brownian-excursion limit distribution.
\end{proposition}

\begin{warning}[Ambient rarity is not an arithmetic restriction]
\label{warn:ambient-rarity}
The estimates in Proposition~\ref{prop:height-transition} describe the uniform Catalan
ensemble.  A deterministic arithmetic selector may choose a member of an exponentially
small Catalan sector.  No conclusion for $\Mert$ follows until an arithmetic property of
the selector is proved.
\end{warning}

A simple reflection estimate gives a finite bound useful for calibration.

\begin{lemma}[Finite high-height tail \DED]
\label{lem:finite-height-tail}
For $1\le H\le N$,
\begin{equation}
\frac{\#\{D\in\mathcal D_N:\max D\ge H\}}{C_N}
\le
(N+1)\exp\!\left(-\frac{H^2}{2N}\right).
\label{eq:finite-height-tail}
\end{equation}
\end{lemma}

\begin{proof}
A Dyck path reaching $H$ is in particular an unrestricted bridge from $0$ to $0$ that
reaches $H$.  Reflection after its first visit to $H$ injects such bridges into walks from
$0$ to $2H$, giving at most $\binom{2N}{N+H}$ possibilities.  Since
$C_N=(N+1)^{-1}\binom{2N}{N}$,
\begin{equation}
\frac{\#\{D:\max D\ge H\}}{C_N}
\le
(N+1)\prod_{j=1}^{H}\frac{N-j+1}{N+j}.
\label{eq:tail-product}
\end{equation}
Using $\log(1-u)\le-u$ on the product yields \eqref{eq:finite-height-tail}.
\end{proof}

\section{Occupation energy and the exact operator \texorpdfstring{$J_H^*J_H$}{J*J}}
\label{sec:occupation-energy}

For $q\ge1$ define the ordinary-time Catalan occupation energy
\begin{equation}
\mathcal E_q(x,H)
=
\sum_{h=1}^{H}|\Mert(x+h)-\Mert(x)|^{2q}.
\label{eq:occupation-energy}
\end{equation}
The compressed analogue is
\begin{equation}
\mathcal E_q^{\#}(m,L)
=
\sum_{u=1}^{L}|X_{m+u}-X_m|^{2q}.
\label{eq:compressed-energy}
\end{equation}
These quantities should not be silently identified; ordinary time repeats heights at nonsquarefree
integers.

\begin{definition}[Dyadic Catalan energy]
\label{def:dyadic-energy}
For fixed $q\ge1$, write $\mathrm{CE}^{\mathrm{dyad}}_q$ for the estimate
\begin{equation}
\mathcal E_q(x,x)\ll_{q,\eps}x^{q+1+\eps}
\quad(x\ge2).
\label{eq:dyadic-energy}
\end{equation}
\end{definition}

\begin{theorem}[Exact exponent ladder \DED]
\label{thm:energy-ladder}
For every fixed $q\ge1$,
\begin{equation}
\mathrm{CE}^{\mathrm{dyad}}_q
\quad\Longrightarrow\quad
\Mert(X)\ll_{q,\eps}X^{(q+1)/(2q+1)+\eps}.
\label{eq:energy-ladder}
\end{equation}
Moreover,
\begin{equation}
\bigl[\mathrm{CE}^{\mathrm{dyad}}_q\text{ for every fixed }q\bigr]
\quad\Longleftrightarrow\quad
\mathrm{RH}.
\label{eq:all-q-rh}
\end{equation}
\end{theorem}

\begin{proof}
Put $S_x(h)=\Mert(x+h)-\Mert(x)$ and $A=|S_x(x)|$.  Since the path is
one-Lipschitz, at least a constant multiple of $A$ final values have size at least a
constant multiple of $A$.  Hence
\begin{equation}
A^{2q+1}\ll_q\mathcal E_q(x,x).
\label{eq:plateau-energy}
\end{equation}
Estimate \eqref{eq:energy-ladder} follows for $\Mert(2x)-\Mert(x)$ and then by dyadic
telescoping.  If the hierarchy holds for all fixed $q$, the exponents decrease to
$1/2$.  Conversely, RH gives $|S_x(h)|\ll_\eta x^{1/2+\eta}$ uniformly for
$x\le x+h\le2x$, which implies \eqref{eq:dyadic-energy} after choosing $\eta$ in
terms of $q$ and $\eps$.
\end{proof}

The quadratic member has an exact spectral structure.  Put
\begin{equation}
a_n=\Mob(x+n),\qquad 1\le n\le H,
\label{eq:increment-vector}
\end{equation}
and let $J_H$ be the cumulative-sum matrix,
\begin{equation}
(J_Ha)_h=\sum_{n\le h}a_n.
\label{eq:summation-matrix}
\end{equation}
Then
\begin{equation}
\mathcal E_1(x,H)=\langle a,K_Ha\rangle,
\qquad K_H=J_H^*J_H,
\label{eq:quadratic-kernel}
\end{equation}
with
\begin{equation}
(K_H)_{rs}=H+1-\max(r,s).
\label{eq:quadratic-kernel-entries}
\end{equation}

\begin{proposition}[Exact occupation spectrum \DED]
\label{prop:occupation-spectrum}
The eigenvalues of $K_H$ are
\begin{equation}
\lambda_k=
\frac{1}{4\sin^2\!\left(\frac{(2k-1)\pi}{4H+2}\right)},
\qquad 1\le k\le H.
\label{eq:occupation-eigenvalues}
\end{equation}
In particular, for $k\ll H$,
\begin{equation}
\lambda_k\asymp\frac{H^2}{(2k-1)^2}.
\label{eq:low-mode-weight}
\end{equation}
\end{proposition}

Thus low Fourier-like modes are weighted much more strongly than generic high modes.  This is the
reason to retain prime colours or twisted sums until after the $J_H^*J_H$ weighting rather than
projecting immediately to a scalar endpoint.

Let
\begin{equation}
A_{x,H}(\alpha)=\sum_{n=1}^{H}\Mob(x+n)e^{2\pi i\alpha n}.
\label{eq:short-twisted-mobius}
\end{equation}
The sine eigenvectors of $K_H$ give the deterministic comparison
\begin{equation}
\boxed{
\mathcal E_1(x,H)<
(H+1)\sup_{\alpha\in\R}|A_{x,H}(\alpha)|^2.}
\label{eq:energy-twist-bound}
\end{equation}

\section{What is already known analytically, and the first power frontier}
\label{sec:known-analytic-frontier}

Matom\"aki and Ter\"av\"ainen proved cancellation of $\Mob$ in every interval of length
$x^\theta$ for every $\theta>0.55$, with Ramar\'e extraction as a key new ingredient
\cite{MatomakiTeravainen2023}.  They also proved a uniform twisted estimate for
$\theta>3/5$; higher-uniformity results in all intervals at the threshold $5/8$ were later
proved by Matom\"aki, Shao, Tao, and Ter\"av\"ainen \cite{MatomakiShaoTaoTeravainen2023}.
Combining the twisted estimates with \eqref{eq:energy-twist-bound} gives logarithmic
anti-ballistic occupation-energy savings in their valid ranges.

\begin{proposition}[Known anti-ballistic quadratic bounds \LIT+\DED]
\label{prop:known-antiballistic}
Fix $\eta>0$.  The published twisted short-interval estimates imply bounds of the form
\begin{equation}
\mathcal E_1(x,H)
\ll_{\theta,\eta}
\frac{H^3}{(\log x)^{2/3-2\eta}}
\qquad(H\ge x^\theta,\ \theta>3/5),
\label{eq:known-log-energy}
\end{equation}
and, for every $A>0$,
\begin{equation}
\mathcal E_1(x,H)
\ll_{A,\eta}
\frac{H^3}{(\log x)^A}
\qquad(x^{5/8+\eta}\le H\le x^{1-\eta}).
\label{eq:known-arbitrary-log-energy}
\end{equation}
\end{proposition}

These are genuine improvements over the trivial ballistic energy $O(H^3)$, but they are not power
savings in $H$.

\begin{lemma}[One-sided tent conversion \DED]
\label{lem:tent}
For any integer-valued path $S(0)=0$ with increments of modulus at most one, let
$\Delta=\max_{0\le h\le H}|S(h)|$.  Then
\begin{equation}
\frac{\Delta^3}{3}
\le
\sum_{h=1}^{H}|S(h)|^2.
\label{eq:tent-inequality}
\end{equation}
\end{lemma}

\begin{proof}
If $|S(h_0)|=\Delta$, one-Lipschitz continuity forces
$|S(h_0-t)|\ge\Delta-t$ for $0\le t\le\Delta$.  Summing the resulting squares gives
\eqref{eq:tent-inequality}.
\end{proof}

\begin{conjecture}[First major open energy target \OPEN]
\label{conj:first-power-saving}
There exist fixed $\theta,\delta>0$ such that, uniformly in the initial point on dyadic ranges and
for $H\asymp X^\theta$,
\begin{equation}
\mathcal E_1(x,H)\ll H^{3-\delta}X^\eps
\label{eq:first-power-saving}
\end{equation}
for every $\eps>0$.
\end{conjecture}

\begin{proposition}[Propagation of a mesoscopic power saving \DED]
\label{prop:power-propagation}
If \eqref{eq:first-power-saving} holds, then
\begin{equation}
\Mert(X)\ll X^{1-\theta\delta/3+\eps}.
\label{eq:power-propagation}
\end{equation}
\end{proposition}

\begin{proof}
Lemma~\ref{lem:tent} gives a window endpoint bound
$O(H^{1-\delta/3}X^\eps)$.  Partition a dyadic interval into $O(X/H)$ windows and sum the
endpoint increments.
\end{proof}

\begin{remark}[Calibration]
\label{rem:energy-calibration}
The Brownian target $\mathcal E_1\ll H^{2+\eps}$ is a long-term endpoint, not the first milestone.
The first nonclassical analytic step is any fixed power saving in \eqref{eq:first-power-saving}.
\end{remark}

\section{Why pointwise revelation is insufficient}
\label{sec:pointwise-no-go}

Suppose a word of length $m$ has a set $F$ of already fixed signs.  Define
\begin{align}
s_F(t)&=\sum_{\substack{j\le t\\j\in F}}\sigma_j,
\label{eq:fixed-prefix-sum}\\
U_F(t)&=\#\{j\le t:j\notin F\}.
\label{eq:unrevealed-count}
\end{align}

\begin{proposition}[Exact pointwise envelope \DED]
\label{prop:pointwise-envelope}
For every prefix $t$,
\begin{equation}
\max_{\text{compatible completions}}|S_t|
=U_F(t)+|s_F(t)|.
\label{eq:pointwise-envelope}
\end{equation}
Consequently,
\begin{equation}
H_F^*(m)
=
\max_{t\le m}\bigl(U_F(t)+|s_F(t)|\bigr).
\label{eq:pointwise-global-envelope}
\end{equation}
\end{proposition}

\begin{proof}
Every free sign in a fixed prefix can be chosen to reinforce the sign of $s_F(t)$; if the fixed sum is
zero, choose all free signs with the same orientation.
\end{proof}

\begin{corollary}[Near-total revelation barrier \NEG]
\label{cor:near-total-revelation}
A pointwise-revelation mechanism can force the envelope to
$O(m^{1/2+\eps})$ only if every relevant prefix contains at most
$O(m^{1/2+\eps})$ unrevealed positions, apart from cancellation already proved in the fixed part.
\end{corollary}

Thus a proof based only on revealing $\omega(n)$ layers one position at a time approaches full knowledge of the
M\"obius word and risks becoming the desired Mertens estimate in disguise.  Collective restrictions are required.

\section{Consistent signed graphs and collective prefix geometry}
\label{sec:signed-graph}

A convenient universal model is a consistent signed graph on positions.  If an edge $uv$ carries
$\eta_{uv}\in\{\pm1\}$, impose
\begin{equation}
\varepsilon_v=\eta_{uv}\varepsilon_u.
\label{eq:signed-graph-relation}
\end{equation}
On each connected component $C$ there are deterministic labels $\lambda_v\in\{\pm1\}$ and one
free orientation $s_C$ such that $\varepsilon_v=s_C\lambda_v$.

\begin{definition}[Prefix charge]
\label{def:prefix-charge}
For a component $C$ define
\begin{equation}
Q_C(t)=\sum_{\substack{v\in C\\v\le t}}\lambda_v.
\label{eq:prefix-charge}
\end{equation}
\end{definition}

\begin{proposition}[Exact signed-graph envelope \DED]
\label{prop:signed-graph-envelope}
For every prefix $t$,
\begin{equation}
\max_{\text{compatible orientations}}|S_t|
=
\sum_C|Q_C(t)|,
\label{eq:signed-graph-envelope}
\end{equation}
and therefore
\begin{equation}
H_G^*=\max_t\sum_C|Q_C(t)|.
\label{eq:signed-graph-global-envelope}
\end{equation}
\end{proposition}

This identity converts an arithmetic relation graph into an exact geometric height problem.  The next sections
specialize it to multiplicative divisor components.

\section{Adversarial multiplicative completions and signed displacement cuts}
\label{sec:extinction}

A separate project route begins with the finite cube of multiplicative squarefree sign completions.  Fix the
squarefree support and prime signs up to scale $N$.  The true M\"obius word is one vertex of this finite cube.
For a completion $x$ define
\begin{equation}
C_h(x)=\sum_{n\le N-h}x_nx_{n+h},
\qquad
d_h(x)=C_h(x)-C_h(\Mob).
\label{eq:completion-correlation}
\end{equation}
Let $D(x)=\{n:x_n\ne\Mob(n)\}$ on the squarefree support.

\begin{theorem}[Exact signed-cut identity \DED]
\label{thm:signed-cut}
For each displacement $h$, there is a signed displacement graph $G_h$ such that
\begin{equation}
\frac{d_h(x)}{2}=\Cut_{G_h}(D(x)).
\label{eq:signed-cut}
\end{equation}
If $u_e\in\Ftwo^{\pi(N)-1}$ is the odd-prime parity form attached to an edge $e$, its sign satisfies
\begin{equation}
w_h(e)=(-1)^{|u_e|+h+1}.
\label{eq:edge-parity-law}
\end{equation}
\end{theorem}

The exact finite formulas convert every displacement test into a weighted parity-threshold condition on the
prime-sign Boolean cube.  Neighboring literature on random multiplicative functions, small partial sums, and sign
patterns includes \cite{AymoneSidoravicius2017,KoukoulopoulosSoundararajan2020,MatomakiRadziwillTao2016}; the
specific finite extinction architecture is treated here only as a project construction, without a literature-wide
priority assertion.

\begin{experiment}[Finite extinction audit \COMP]
\label{exp:finite-extinction}
The project computations verified complete extinction of every non-M\"obius multiplicative completion at several
small finite scales, including $N=25,43,53,61,71,73$.  This is finite evidence only.  The general statement
\begin{equation}
\mathcal B_{H_0(N)}(N)=\{0\}
\label{eq:general-extinction-target}
\end{equation}
is not proved.
\end{experiment}

\begin{warning}[Extinction route versus mean square]
\label{warn:extinction-nonroute}
A Fej\'er or moving-window mean-square estimate can be formulated directly for the true $\Mob$ sequence without the
completion cube.  Such an estimate may be mathematically useful, but it does not explain finite extinction and is not
an internal continuation of that mechanism.
\end{warning}

\section{Collective divisor-cube capacity}
\label{sec:divisor-capacity}

We now introduce a collective arithmetic restriction that acts on whole multiplicative components rather than fixing
signs position by position.  Let
\begin{equation}
P=\prod_{p\in\mathcal P}p
\label{eq:finite-prime-product}
\end{equation}
be squarefree.  Every squarefree integer has a unique representation
\begin{equation}
n=ad,\qquad d\mid P,\qquad (a,P)=1.
\label{eq:component-factorization}
\end{equation}
The relations generated by the primes dividing $P$ force every compatible sign assignment to have the form
\begin{equation}
\varepsilon(ad)=s_a\Mob(d),\qquad s_a\in\{\pm1\}.
\label{eq:component-signs}
\end{equation}
Thus one bit remains for each $P$-free squarefree base $a$.

Fix a window $W=(A,T]$ and define
\begin{equation}
D_a(W)=\{d\mid P:A<ad\le T\}.
\label{eq:active-divisors}
\end{equation}
Let
\begin{align}
E_a(W)&=\#\{d\in D_a(W):\Mob(d)=+1\},
\label{eq:even-parity-count}\\
O_a(W)&=\#\{d\in D_a(W):\Mob(d)=-1\}.
\label{eq:odd-parity-count}
\end{align}

\begin{theorem}[Exact component capacity \DED]
\label{thm:component-capacity}
Within the component $S_a=\{ad:d\in D_a(W)\}$ the number of up-steps satisfies
\begin{equation}
|B\cap S_a|\le\max(E_a,O_a).
\label{eq:component-up-capacity}
\end{equation}
Moreover, if only the multiplicative relations from $P$ are imposed, then the exact maximum signed contribution over
the whole window is
\begin{equation}
H_P^{\mathrm{free}}(A,T)
=
L(A,T)-2\mathcal B_P(A,T),
\label{eq:component-exact-height}
\end{equation}
where
\begin{align}
L(A,T)&=\Sq(T)-\Sq(A),
\label{eq:window-squarefree-length}\\
\mathcal B_P(A,T)&=\sum_a\min(E_a,O_a).
\label{eq:capacity-deficit}
\end{align}
\end{theorem}

\begin{proof}
In the component $a$, all signs equal $s_a\Mob(d)$.  If $s_a=+1$, the up-steps are the $E_a$ divisors; if
$s_a=-1$, they are the $O_a$ divisors.  Hence the maximal signed contribution of that component is
$|E_a-O_a|=(E_a+O_a)-2\min(E_a,O_a)$.  Sum over the disjoint components.
\end{proof}

\begin{remark}[Sign-sensitive fail gate]
\label{rem:capacity-sign-sensitive}
The deficit in \eqref{eq:capacity-deficit} genuinely uses M\"obius parity.  Under the replacement
$\Mob\mapsto|\Mob|$, all active divisors in a component have the same sign and the deficit disappears.  Thus this
capacity passes the sign-sensitivity test required of a nontrivial arithmetic layer.
\end{remark}

For a multiplicative window $T=\lambda A$, define
\begin{equation}
g_{P,\lambda}(y)
=
\sum_{\substack{d\mid P\\y<d\le\lambda y}}\Mob(d).
\label{eq:divisor-window-kernel}
\end{equation}

\begin{theorem}[Exact asymptotic capacity ratio for fixed $P$ \DED]
\label{thm:rho-capacity}
For fixed squarefree $P$ and $\lambda>1$,
\begin{equation}
H_P(A,\lambda A)
=
\sum_{\substack{a\ \mathrm{squarefree}\\(a,P)=1}}
\left|g_{P,\lambda}(A/a)\right|.
\label{eq:capacity-sum-kernel}
\end{equation}
Furthermore,
\begin{equation}
H_P(A,\lambda A)
=
\delta_PA I_P(\lambda)+O_P(\sqrt A),
\label{eq:capacity-asymptotic}
\end{equation}
where
\begin{align}
\delta_P&=\frac6{\pi^2}\prod_{p\mid P}\frac{p}{p+1},
\label{eq:coprime-squarefree-density}\\
I_P(\lambda)&=
\int_{1/\lambda}^{\infty}
\left|
\sum_{\substack{d\mid P\\y<d\le\lambda y}}\Mob(d)
\right|\frac{dy}{y^2}.
\label{eq:capacity-integral}
\end{align}
Consequently the surviving ballistic ratio is
\begin{equation}
\boxed{
\rho_P(\lambda)
=
\frac{1}{\lambda-1}
\left(\prod_{p\mid P}\frac{p}{p+1}\right)
I_P(\lambda).}
\label{eq:rho-definition}
\end{equation}
\end{theorem}

\begin{proof}
For each base $a$, the signed component sum is $s_ag_{P,\lambda}(A/a)$, proving
\eqref{eq:capacity-sum-kernel}.  The count of squarefree $a$ coprime to fixed $P$ has density $\delta_P$ and an
$O_P(\sqrt A)$ elementary squarefree-counting error.  Since $g_{P,\lambda}$ is a finite step function, summation over
its finitely many threshold intervals gives \eqref{eq:capacity-asymptotic}.  Divide by
$\Sq(\lambda A)-\Sq(A)=(6/\pi^2)(\lambda-1)A+O(\sqrt A)$.
\end{proof}

\begin{example}
\label{ex:p30-capacity}
For $P=30$ and $\lambda=2$, the opposite-parity divisors $5$ and $6$ already force the full deficit and one obtains
\begin{equation}
\rho_{30}(2)=\frac89.
\label{eq:rho30}
\end{equation}
For $P=210$ the corresponding exact finite-divisor calculation gives
\begin{equation}
\rho_{210}(2)=\frac56.
\label{eq:rho210}
\end{equation}
These are constant-factor contractions, not exponent contractions.
\end{example}

\section{Prime insertion, monotonicity, and the entropy--capacity barrier}
\label{sec:capacity-monotonicity}

Set
\begin{equation}
F_{P,\lambda}(x)
=
\sum_{d\mid P}\Mob(d)\mathbf 1_{(1/d,\lambda/d]}(x)
\label{eq:reciprocal-kernel}
\end{equation}
and
\begin{equation}
c(P)=\prod_{p\mid P}\frac{p}{p+1}.
\label{eq:cP}
\end{equation}
Then
\begin{equation}
\rho_P(\lambda)
=
\frac{c(P)}{\lambda-1}\|F_{P,\lambda}\|_{L^1(0,\infty)}.
\label{eq:rho-L1}
\end{equation}

\begin{proposition}[Delayed anti-copy under prime insertion \DED]
\label{prop:prime-anti-copy}
If $p\nmid P$, then
\begin{equation}
F_{Pp,\lambda}(x)
=
F_{P,\lambda}(x)-F_{P,\lambda}(px).
\label{eq:prime-anti-copy}
\end{equation}
Consequently,
\begin{equation}
\frac{p-1}{p+1}\rho_P(\lambda)
\le
\rho_{Pp}(\lambda)
\le
\rho_P(\lambda).
\label{eq:rho-two-sided-step}
\end{equation}
\end{proposition}

\begin{proof}
Split the divisors of $Pp$ into $d$ and $pd$ with $d\mid P$, using $\Mob(pd)=-\Mob(d)$.
For the upper bound, apply the triangle inequality to \eqref{eq:prime-anti-copy}; for the lower bound apply the reverse
triangle inequality and use
$\|F_{P,\lambda}(p\,\cdot)\|_1=p^{-1}\|F_{P,\lambda}\|_1$.  Multiplication by
$c(Pp)=p(p+1)^{-1}c(P)$ gives \eqref{eq:rho-two-sided-step}.
\end{proof}

\begin{corollary}[Structural lower barrier \DED]
\label{cor:rho-barrier}
For every finite squarefree $P$,
\begin{equation}
\prod_{p\mid P}\frac{p-1}{p+1}
\le
\rho_P(\lambda)
\le1.
\label{eq:rho-product-barrier}
\end{equation}
In particular,
\begin{equation}
\rho_P(\lambda)
\ge
\frac{c(P)^2}{\zeta(2)}.
\label{eq:rho-zeta-barrier}
\end{equation}
\end{corollary}

\begin{proof}
Iterate \eqref{eq:rho-two-sided-step} from $P=1$.  Since
\begin{equation}
\prod_{p\mid P}\frac{p-1}{p+1}
=c(P)^2\prod_{p\mid P}\left(1-\frac1{p^2}\right),
\label{eq:barrier-factorization}
\end{equation}
and the finite Euler product is at least $1/\zeta(2)$, \eqref{eq:rho-zeta-barrier} follows.
\end{proof}

The exact gain from a newly inserted prime can also be isolated.  Define
\begin{equation}
J_{P,p}(\lambda)
=
\int_{\{F_{P,\lambda}(x)F_{P,\lambda}(px)>0\}}
\min\{|F_{P,\lambda}(x)|,|F_{P,\lambda}(px)|\}\,dx.
\label{eq:overlap-functional}
\end{equation}
Then the identity $|u-v|=|u|+|v|-2\min(|u|,|v|)$ on the same-sign region gives
\begin{equation}
\rho_P(\lambda)-\rho_{Pp}(\lambda)
=
\frac{2p}{p+1}\frac{c(P)}{\lambda-1}J_{P,p}(\lambda).
\label{eq:rho-synergy}
\end{equation}
Thus the layer synergy is exactly the same-sign overlap between a divisor kernel and its multiplicative dilation.

For the primorial $P_z=\prod_{p\le z}p$, Mertens' product theorem gives
\begin{equation}
c(P_z)\sim\frac{\zeta(2)e^{-\gamma}}{\log z}.
\label{eq:c-primorial}
\end{equation}
Hence \eqref{eq:rho-zeta-barrier} implies the structural lower scale
\begin{equation}
\rho_{P_z}(\lambda)
\gtrsim
\frac{\zeta(2)e^{-2\gamma}}{(\log z)^2}.
\label{eq:rho-primorial-lower}
\end{equation}

\begin{warning}[Entropy--capacity tradeoff \NEG]
\label{warn:entropy-capacity}
Driving $\rho_{P_z}$ toward zero by taking $z\to\infty$ also drives the density of $P_z$-free base components toward
zero.  The number of independent orientations $s_a$ in a macroscopic interval is of order $A/\log z$ in the usual
sieve range, so the resulting number of componentwise free completions is only
\begin{equation}
2^{O(A/\log z)}=2^{o(A)}
\qquad(z\to\infty).
\label{eq:entropy-loss}
\end{equation}
Thus this within-component mechanism cannot simultaneously produce sublinear confinement and retain
$2^{cA}$ independent component orientations with fixed $c>0$.
\end{warning}

\section{Static Ramar\'e does not couple divisor components}
\label{sec:static-ramare-no-go}

The Catalan programme uses the exact Ramar\'e colouring
\begin{equation}
\Mob(n)
=
-\sum_{\substack{p\in\mathcal P\\p\mid n}}
\frac{\Mob(n/p)}{\nu_{\mathcal P}(n)},
\qquad
\nu_{\mathcal P}(n)=\#\{p\in\mathcal P:p\mid n\},
\label{eq:ramare-colouring}
\end{equation}
for squarefree $n$ having at least one prime factor in the band $\mathcal P$.  This weighted extraction is a central
analytic device in the short-interval work of Matom\"aki and Ter\"av\"ainen \cite{MatomakiTeravainen2023}.

\begin{proposition}[Static component decomposition \NEG]
\label{prop:static-ramare-components}
Fix a finite prime set $\mathcal P$ and let $P=\prod_{p\in\mathcal P}p$.  Form the graph on squarefree integers with an
edge between $n$ and $n/p$ whenever $p\in\mathcal P$ divides $n$.  Its connected components are exactly
\begin{equation}
\mathcal C_a=\{ad:d\mid P\},
\qquad (a,P)=1,
\label{eq:ramare-components}
\end{equation}
with the obvious truncation if a finite cutoff is imposed.  Every sign assignment compatible with all static prime-deletion
relations has one free bit per component and has the form \eqref{eq:component-signs}.
\end{proposition}

\begin{proof}
Deleting primes from $\mathcal P$ changes only the divisor $d$ and leaves the $P$-free part $a$ unchanged.  Conversely,
any two divisors of $P$ can be connected by successive insertion or deletion of primes from $\mathcal P$.  M\"obius parity
changes once per prime insertion, yielding \eqref{eq:component-signs}.
\end{proof}

\begin{corollary}[No static inter-component gain \NEG]
\label{cor:no-static-intercomponent}
A reweighting of the static identities \eqref{eq:ramare-colouring} alone cannot correlate the independent orientations
$s_a$ of distinct components.  Any gain beyond Theorem~\ref{thm:rho-capacity} must use additional information, such as the
causal order/cutoff, an analytic recombination estimate, or another relation coupling different bases.
\end{corollary}

This no-go is compatible with the ``redundant prime lift'' already present in the Catalan notebook: subtracting the exact
prime-shell identity at two excursion endpoints reproduces the same Mertens prefix for every prime and gives no estimate by
averaging over primes.

\section{The causal Ramar\'e prefix lift and the open inter-component frontier}
\label{sec:causal-ramare}

For an interval beginning at $a$, retain the prime colour until after the prefix has been formed.  Define the coloured currents
\begin{equation}
C_p(u)
=
-\sum_{\substack{a<n\le a+u\\p\mid n,\ \Mob(n)^2=1\\\nu_{\mathcal P}(n)>0}}
\frac{\Mob(n/p)}{\nu_{\mathcal P}(n)}
\label{eq:coloured-current}
\end{equation}
and the rough remainder
\begin{equation}
R_{\mathcal P}(u)
=
\sum_{\substack{a<n\le a+u\\\nu_{\mathcal P}(n)=0}}\Mob(n).
\label{eq:rough-current}
\end{equation}
Then every prefix satisfies the exact vector identity
\begin{equation}
\boxed{
\Mert(a+u)-\Mert(a)
=
R_{\mathcal P}(u)+\sum_{p\in\mathcal P}C_p(u).}
\label{eq:ramare-prefix-vector}
\end{equation}
Unlike the static component relations, the functions $u\mapsto C_p(u)$ retain the causal truncation and can be passed through
the occupation operator $J_H^*J_H$ before scalar projection.

\begin{conjecture}[Weak causal prime-incidence power saving \OPEN]
\label{conj:weak-causal-frame}
There exist admissible growing prime bands and fixed $\delta>0$ such that the signed recombination in
\eqref{eq:ramare-prefix-vector}, after low-mode $J_H^*J_H$ weighting and with the rough term retained, implies
\begin{equation}
\mathcal E_1(x,H)\ll H^{3-\delta}X^\eps
\label{eq:causal-power-target}
\end{equation}
uniformly on a nontrivial polynomial window range.
\end{conjecture}

\begin{remark}[Long-term Brownian frame]
\label{rem:brownian-frame}
The stronger square-function target of order $H^{2+\eps}$ is retained as the Brownian endpoint of the programme, not as the
first milestone.  The finite-power estimate \eqref{eq:causal-power-target} is the calibrated first open target.
\end{remark}

\section{Catalan rank and collective capacity constraints}
\label{sec:catalan-rank}

A generic collective restriction on up-step positions has the form
\begin{equation}
|B\cap S|\le c(S).
\label{eq:generic-capacity}
\end{equation}
In the Catalan matroid description, a Dyck basis with ordered up-step positions $b_1<\cdots<b_n$ satisfies
\begin{equation}
b_i\le2i-1
\qquad(1\le i\le n).
\label{eq:catalan-basis-criterion}
\end{equation}
The project notebook develops a universal rank-to-height conversion for restrictions of this type.

\begin{proposition}[Exact layer drop in a Catalan minor \DED]
\label{prop:layer-drop}
Suppose a Catalan completion class has current free prefix set $E_t$ and matroid rank $r_M(E_t)$.  A new layer fixes previously
free sets $P_+$ upward and $P_-$ downward, producing a new minor $M'$.  Define
\begin{equation}
\Delta_t
=
r_M(E_t)-|P_+\cap I_t|
-r_{M'}\bigl(E_t\setminus(P_+\cup P_-)\bigr).
\label{eq:rank-deficit}
\end{equation}
Then $\Delta_t\ge0$ and the maximal compatible prefix height drops by
\begin{equation}
h^*_{\mathrm{old}}(t)-h^*_{\mathrm{new}}(t)=2\Delta_t.
\label{eq:exact-layer-drop}
\end{equation}
\end{proposition}

The divisor-cube capacity in Section~\ref{sec:divisor-capacity} supplies an explicit arithmetic source of a collective deficit;
the unresolved task is to obtain a growing family whose joint deficit acts at mesoscopic scales without spending essentially all
remaining entropy.

\section{Sign constraints and exact ballot suppression}
\label{sec:sign-constrained-words}

The project chat on Mertens sign changes also gives a useful negative calibration.  Fix the actual zero positions of $\Mob$ and
require a candidate ternary word to have partial sums with the same sign pattern as $\Mert$.  In the initial example through
$n=20$, after the forced touch at zero and the first negative steps, the remaining problem is a one-sided ballot problem.

\begin{proposition}[Closed count for the initial sign-constrained model \DED]
\label{prop:sign-model-count}
Let $n$ denote the number of nonzero steps in the model described above.  For $n\ge4$, the number of admissible words is
\begin{equation}
R(n)=\binom{n-3}{\lfloor(n-3)/2\rfloor}.
\label{eq:sign-model-count}
\end{equation}
Thus, among the $2^n$ unrestricted sign words,
\begin{equation}
\frac{R(n)}{2^n}
\sim
\frac18\sqrt{\frac{2}{\pi(n-3)}}.
\label{eq:sign-model-fraction}
\end{equation}
\end{proposition}

\begin{proof}
After the finite forced prefix, reflect the negative walk to a walk starting at height two and require it to avoid zero.  The
reflection principle gives a difference of binomial coefficients for each endpoint.  Summing over admissible endpoints telescopes
to the central binomial expression in \eqref{eq:sign-model-count}.  Stirling's formula gives
\eqref{eq:sign-model-fraction}.
\end{proof}

\begin{remark}[Interpretation \NEG]
\label{rem:sign-model-negative}
In this model, sign/crossing information suppresses the ambient word count only by a polynomial factor.  It is standard ballot
geometry, not yet a M\"obius-specific cancellation mechanism.
\end{remark}

\section{Ordinal-prime transform: exact information and coercive negative excursions}
\label{sec:ordinal-prime}

A complementary project assigns ordinal primes separately to positive and negative nonzero M\"obius occurrences, producing a
sequence $S(n)$ and partial sum $SS(n)$.  Its endpoint theory is exact but largely determined by the pair $(\Mert(n),\Sq(n))$.
After subtracting the deterministic zero contribution, the remaining term is a signed block of $|\Mert(n)|$ consecutive primes.
Thus its endpoint RH criteria are reformulations rather than independent constraints.

The transform nevertheless contains a useful pathwise asymmetry on negative excursions.  Let
\begin{equation}
R(q)=\Mert(q_q)
\label{eq:ordinal-squarefree-walk}
\end{equation}
be the compressed walk in the notation of the ordinal project.  A complete negative excursion has depth
\begin{equation}
H=-\min R(q).
\label{eq:negative-depth}
\end{equation}
For each level $-k$, let $u_k,v_k$ be its first and last visits and write $v_k-u_k=2r_k$.

\begin{theorem}[Coercive negative-excursion energy \DED]
\label{thm:ordinal-cubic-energy}
The ordinal-prime level-return functional $\mathcal A$ defined by the positive prime-separation sums satisfies
\begin{equation}
\mathcal A\ge\frac{H^3-H}{3}.
\label{eq:ordinal-cubic-energy}
\end{equation}
More sharply, if
\begin{equation}
G(k)=\inf_{m\ge2}(p_{m+k}-p_m),
\label{eq:prime-gap-block-min}
\end{equation}
then
\begin{equation}
\mathcal A\ge\sum_{k=1}^{H-1}(H-k)G(k).
\label{eq:ordinal-sharp-energy}
\end{equation}
\end{theorem}

\begin{warning}[What the cubic lower bound does not do]
\label{warn:ordinal-upper-missing}
Theorem~\ref{thm:ordinal-cubic-energy} becomes an upper bound on the excursion depth only after an independent upper estimate for
$\mathcal A$ is proved.  No such estimate at the required scale is presently supplied by the transform itself.
\end{warning}

\section{The analytic \texorpdfstring{$\alpha$}{alpha} programme as a calibrated companion}
\label{sec:alpha-programme}

The ordinal programme also produces Dirichlet-type functions whose convergence boundaries encode the zeta zero frontier.  Let
\begin{equation}
\Theta=\sup_{\zeta(\rho)=0}\Re\rho.
\label{eq:theta-zero-frontier}
\end{equation}
The source notebook proves, for its explicitly defined $\alpha$ and paired term $\delta$, the abscissa identity
\begin{equation}
\sigma_c(\alpha)=\sigma_c(\delta)=1+\Theta.
\label{eq:alpha-abscissa}
\end{equation}
It also isolates a component $B(s)$ whose poles occur at $1+\rho$ with multiplicity inherited from the zeros of $\zeta$.

\begin{warning}[Pole transfer boundary]
\label{warn:alpha-poles}
The poles of the explicit $B$ component do not by themselves prove corresponding poles for the full $\alpha$ or $\delta$.
A relatively smaller Mellin remainder can still generate a logarithmic branch singularity.  Continuation of the full object requires
an independent uniform estimate.
\end{warning}

An exact energy version is particularly clean.  If
\begin{equation}
\mathcal I(\sigma)
=
\int_1^\infty
\bigl(SS(x)-\lfloor x\rfloor+\Sq(x)\bigr)^2
x^{-2\sigma-1}\,dx,
\label{eq:alpha-energy}
\end{equation}
then the project establishes an exact comparison with a weighted Mertens mean square and consequently
\begin{equation}
\mathrm{RH}
\quad\Longleftrightarrow\quad
\mathcal I(\sigma)<\infty
\text{ for every }\sigma>3/2.
\label{eq:alpha-energy-rh}
\end{equation}
This is a useful calibration of exponent strength, not an independent proof mechanism.

\section{Divisibility surfaces and \texorpdfstring{$p$-adic}{p-adic} shell coordinates}
\label{sec:surface-padic}

The divisibility-surface programme introduces
\begin{equation}
F(n,r)=\sum_{\substack{a\mid n\\a>r}}\Mob(n/a)
\label{eq:divisibility-surface}
\end{equation}
and a finite square matrix $B_N$ built from these values.  One of its exact structural identities is the explicit inverse
\begin{equation}
(B_N^{-1})_{r,n}=\mathbf 1_{n\mid r+1}-\mathbf 1_{n\mid r}.
\label{eq:surface-inverse}
\end{equation}
This supplies a common coordinate system for M\"obius, Liouville, squarefree projections, and additive shifts.  The construction is
structural: by itself it does not create the missing cancellation needed for ordinary fixed-shift Chowla or
\eqref{eq:mertens-rh-criterion}.

A related $p$-adic shell factorization writes suitable Liouville correlations as a valuation-parity factor times a $p$-free part.
The special triadic centering can eliminate a constant local channel, but the remaining additively modulated $p$-free channels still
require cancellation.  This yields another general methodological rule: local martingale orthogonality is not arithmetic
independence.

\begin{warning}[Single-prime shell no-go]
\label{warn:single-prime-no-go}
The single-prime shell route does not provide a universal base-independent contraction.  Exact shell decompositions are retained as
coordinates, but closure of their tails requires cancellation of essentially the original Chowla/Mertens type.  They are not used here
as a proof shortcut.
\end{warning}

\section{One-way Catalan calibration of hypothetical off-critical zeros}
\label{sec:offcritical-calibration}

This section records a consequence, not a proof strategy.  Suppose $\zeta$ has a zero
$\rho=\beta+i\gamma$ with $\beta>1/2$.  The Mellin representation of $1/\zeta$ implies that for every $\eps>0$,
$\Mert(x)$ cannot be $O(x^{\beta-\eps})$; see \cite{Titchmarsh1986}.  Thus there are arbitrarily large points with Mertens displacement at
that exponent scale.

Take such a compressed point $m$, let $z<m$ be the last return to zero, reflect the open meander if necessary, and close it by adding
$|X_m|$ downward steps.  This gives a canonical Dyck path of some semilength $N_m\le m$ and height at least $|X_m|$.  Combining this
with Lemma~\ref{lem:finite-height-tail} yields the following calibration.

\begin{proposition}[Off-critical zero forces rare Catalan selections \DED]
\label{prop:offcritical-catalan}
If $\zeta(\beta+i\gamma)=0$ with $\beta>1/2$, then for every $\eps>0$ there are arbitrarily large compressed times for which the canonical
closed Mertens meander belongs to a Catalan height sector of relative size at most
\begin{equation}
\exp\!\left(-m^{2\beta-1-2\eps+o(1)}\right).
\label{eq:offcritical-catalan-rarity}
\end{equation}
\end{proposition}

\begin{warning}[Not a selector theorem]
\label{warn:offcritical-not-proof}
Proposition~\ref{prop:offcritical-catalan} is one-way.  It says what an off-critical zero would force the arithmetic selector to do.
Excluding all such rare sectors is essentially the unresolved Mertens/RH problem unless exclusion is obtained from an independent
arithmetic mechanism.  No low-height anomaly is inferred from the reflected zero $1-\beta$.
\end{warning}

The functional equation pairs $\beta$ with $1-\beta$, while $1/2$ is its unique fixed point.  Independently, $1/2$ is the Catalan
height-transition exponent of Proposition~\ref{prop:height-transition}.  This coincidence is a useful structural calibration, not an
identity between zeta zeros and Catalan species.

\section{Project computations on the full M\"obius data set}
\label{sec:project-computations}

The following experiments belong to the project data analysis.  They are finite statements and are not used as asymptotic premises.

\begin{experiment}[Excursion census through $455{,}052{,}511$ \COMP]
\label{exp:excursion-census}
A blockwise scan of the binary M\"obius data through
\begin{equation}
N=455{,}052{,}511
\label{eq:data-N}
\end{equation}
reported the following values:
\begin{center}
\begin{tabular}{lr}
\toprule
quantity & measured value\\
\midrule
nonzero squarefree steps & $276{,}638{,}778$\\
$\Mert(N)$ & $-5414$\\
returns of the compressed walk to zero & $53{,}078$\\
real sign changes between excursions & $26{,}720$\\
last ordinary-time return & $420{,}059{,}551$\\
last compressed-time return & $255{,}365{,}638$\\
open compressed tail & $21{,}273{,}140$\\
mean completed-excursion gap & $4811.139041$\\
maximum completed-excursion gap & $33{,}409{,}232$\\
\bottomrule
\end{tabular}
\end{center}
The same audit computed a cumulative Catalan-compatible path count with
$\log_{10}$ fraction approximately $-77{,}335.1889795$ relative to all unconstrained nonzero sign words.  This enormous ambient entropy
loss does not by itself imply height confinement, by Warning~\ref{warn:ambient-rarity}.
\end{experiment}

\begin{experiment}[Finite square-root envelope audit \COMP]
\label{exp:finite-sqrt-envelope}
An exhaustive scan over every integer
\begin{equation}
1000\le x\le455{,}052{,}511
\label{eq:finite-envelope-range}
\end{equation}
found
\begin{equation}
|\Mert(x)|\le0.472203\sqrt x
\label{eq:finite-envelope}
\end{equation}
throughout that finite range, with the maximum normalized value attained at $x=2803$.  On the same squarefree support, $25$ random-sign
controls had mean normalized maximum $1.8972$ and standard deviation $0.3384$; none was below the M\"obius value.
\end{experiment}

\begin{warning}[Finite envelope is not an $O(\sqrt x)$ theorem]
\label{warn:finite-envelope}
Experiment~\ref{exp:finite-sqrt-envelope} is a finite computation only.  It neither proves nor suggests that the false Mertens conjecture
should hold asymptotically.  The correct RH-scale benchmark remains \eqref{eq:mertens-rh-criterion}.
\end{warning}

\begin{experiment}[Catalan height-percentile audit \COMP]
\label{exp:catalan-percentile}
For every completed excursion, the project code compares the observed maximum with all primitive Catalan excursions having the same
compressed length and fixed excursion sign.  Because the height distribution is discrete, the calibrated null statistic is the mid-rank
\begin{equation}
u_r
=
\mathbb P(H<H_r)+\frac12\mathbb P(H=H_r),
\label{eq:catalan-midrank}
\end{equation}
whose expectation under the uniform Catalan null is exactly $1/2$.  An earlier raw geometric average of cumulative probabilities was about
$79.3\%$, but that raw statistic is not uniform under the discrete null and is retained only as a diagnostic.  The mid-rank, not the raw
geometric percentile, is the correct calibrated observable.
\end{experiment}

\section{A unified master spectrum for partial arithmetic information}
\label{sec:master-spectrum}

Let $\mathcal P_r(m)$ be a class of nearest-neighbour paths compatible with the arithmetic information revealed through layer $r$, together
with any explicitly imposed return, sign, endpoint, or meander constraints.  Define
\begin{equation}
\mathfrak Z_{r,m}(v,z,u)
=
\sum_{Y\in\mathcal P_r(m)}
v^{\mathcal R_m(Y)}z^{Y_m}u^{\max_{0\le j\le m}|Y_j|},
\label{eq:master-spectrum}
\end{equation}
where $\mathcal R_m(Y)$ counts returns to zero.

This single polynomial retains three quantities that total completion counts lose: return structure, endpoint displacement, and maximum
height.  The corresponding high-sector count is
\begin{equation}
N_r^{\ge H}(m)
=
\#\{Y\in\mathcal P_r(m):\max_{j\le m}|Y_j|\ge H\}.
\label{eq:high-sector-count}
\end{equation}

\begin{remark}[What counts as genuine progress]
\label{rem:genuine-progress}
A reduction in $|\mathcal P_r(m)|$ is not sufficient.  A useful partial theorem must preferentially reduce
$N_r^{\ge H}(m)$ or an occupation-energy analogue while leaving a nontrivial family of completions alive.  This criterion prevents both
full-revelation tautologies and entropy-only claims.
\end{remark}

\section{Priority research programme}
\label{sec:priority-programme}

The preceding exact results narrow the next steps considerably.

\begin{question}[Causal inter-component capacity]
\label{question:causal-capacity}
Can the causal prime-incidence lift \eqref{eq:ramare-prefix-vector} produce a proved family of constraints
\begin{equation}
|B\cap S|\le c(S)
\label{eq:causal-capacity-question}
\end{equation}
that couples distinct divisor bases $a$ and yields a power-size cumulative prefix deficit before near-total revelation?
\end{question}

\begin{question}[Low-mode Ramar\'e frame]
\label{question:low-mode-frame}
Can one bound the low eigenmodes of $J_H^*J_H$ after Ramar\'e decomposition without summing absolute values of all cross-colour
correlations, and obtain \eqref{eq:causal-power-target}?
\end{question}

\begin{question}[Extinction-to-height transfer]
\label{question:extinction-height}
Can first-witness signed-cut extinction cells be shown to preferentially remove high Catalan sectors, rather than merely reducing the total
number of multiplicative completions?
\end{question}

\begin{question}[Ordinal cubic energy upper bound]
\label{question:ordinal-upper}
Is there an independent upper bound on the negative-excursion functional $\mathcal A$ strong enough, together with
\eqref{eq:ordinal-cubic-energy}, to yield a nontrivial depth estimate?
\end{question}

\begin{question}[Filtered shadow restoration]
\label{question:filtered-shadow}
Bj\"orner's squarefree simplicial complex has Mertens as an Euler-characteristic quantity \cite{Bjorner2011}.  Can a filtered shadow or
persistent-homology inequality force opposite-parity mass to enter within a bounded future multiplicative window, rather than merely
exist somewhere in the static complex?
\end{question}

The priority order is therefore
\begin{equation}
\boxed{
\begin{aligned}
\text{causal arithmetic relation}
&\Longrightarrow \text{ collective prefix/low-mode deficit}\\
&\Longrightarrow \mathcal E_1\ll H^{3-\delta}\\
&\Longrightarrow \text{ global exponent improvement}.
\end{aligned}}
\label{eq:priority-chain}
\end{equation}
The Brownian $H^{2+\eps}$ scale and the all-$q$ RH-complete hierarchy remain later stages.

\section{Negative results and circularity audit}
\label{sec:negative-audit}

The following barriers are retained because they prevent the programme from cycling through equivalent formulations.

\begin{proposition}[Consolidated no-go list \NEG]
\label{prop:no-go-list}
The present collection establishes or records the following limitations.
\begin{enumerate}[label=(\roman*)]
\item Ambient Catalan rarity does not restrict a deterministic arithmetic selector.
\item Pointwise revelation cannot force the RH scale while a macroscopic number of independent signs remain free.
\item Any fixed finite set of multiplicative factor relations leaves a positive linear local capacity in sufficiently short relative windows.
\item Static finite-band Ramar\'e relations have exactly the same $aD(P)$ connected components as the divisor-cube capacity and do not couple their free orientations.
\item A moving prime cutoff can reduce the capacity ratio, but the componentwise free entropy simultaneously loses positive density.
\item Scalar short-interval chord flatness alone does not rule out gentle high mountains; occupation energy contains strictly more path information.
\item The ordinal endpoint transform is a deterministic function of $(\Mert,\Sq)$ and supplies no independent endpoint constraint.
\item The explicit poles of a component of the $\alpha$ decomposition do not automatically transfer to the full function without a continuation estimate.
\item The divisibility surface and $p$-adic shell coordinates are exact changes of structure but do not themselves create a spectral or cancellation gap.
\item The single-prime shell route is blocked unless one proves cancellation in the remaining $p$-free modulated channels.
\end{enumerate}
\end{proposition}

\section{Conclusions}
\label{sec:conclusions}

The most useful common object across the project is not a single transform but the squarefree Mertens path equipped with three compatible
structures: its Catalan return geometry, its multiplicative completion space, and its occupation-energy operator.  The Catalan state space is
exact; the arithmetic selector is the missing mechanism.  The divisor-cube theorem supplies a first explicit M\"obius-specific collective
capacity and quantifies exactly how finite prime layers lower the ballistic constant.  Its lower barrier and entropy tradeoff show why the same
mechanism cannot simply be iterated to the RH scale.  The static Ramar\'e no-go then identifies where genuinely new information must enter:
through causal truncation and analytic recombination of prime-coloured currents, or through another inter-component relation.

The calibrated first analytic objective is therefore not RH and not the Brownian endpoint.  It is the fixed-power occupation saving
\eqref{eq:first-power-saving}.  The calibrated first combinatorial objective is a collective capacity or extinction statement that preferentially
removes high sectors while retaining substantial freedom.  These two objectives are close enough to interact through the low modes of
$J_H^*J_H$, yet weak enough not to be merely restatements of \eqref{eq:mertens-rh-criterion}.

\section*{Notebook history}
\addcontentsline{toc}{section}{Notebook history}

\noindent
\textbf{Version 1, 15 September 2026.}
Consolidated the Catalan--Mertens excursion programme, the occupation-energy/$J^*J$ route, adversarial completions and signed-cut extinction,
the ordinal-prime and $\alpha$ programmes, divisibility-surface/$p$-adic coordinates, the project excursion and sign-change experiments, and
the new divisor-cube collective-capacity analysis.  Added the monotone capacity ratio $\rho_P$, the entropy--capacity barrier, and the static
Ramar\'e component no-go.  No literature-wide novelty claim is made for these notebook deductions.

\begin{thebibliography}{99}

\bibitem{Apostol1976}
T. M. Apostol,
\emph{Introduction to Analytic Number Theory},
Springer-Verlag, New York, 1976.

\bibitem{AymoneSidoravicius2017}
M. Aymone and V. Sidoravicius,
Partial sums of biased random multiplicative functions,
\emph{J. Number Theory} \textbf{172} (2017), 343--382.
\href{https://doi.org/10.1016/j.jnt.2016.08.020}{doi:10.1016/j.jnt.2016.08.020};
arXiv:1412.1157.

\bibitem{Bhat1991}
B. V. Rajarama Bhat,
On greatest common divisor matrices and their applications,
\emph{Linear Algebra Appl.} \textbf{158} (1991), 77--97.
\href{https://doi.org/10.1016/0024-3795(91)90051-W}{doi:10.1016/0024-3795(91)90051-W}.

\bibitem{Bjorner2011}
A. Bj\"orner,
A cell complex in number theory,
\emph{Adv. Appl. Math.} \textbf{46} (2011), 71--85.
\href{https://doi.org/10.1016/j.aam.2010.09.007}{doi:10.1016/j.aam.2010.09.007};
arXiv:1101.5704.

\bibitem{CsakiMohanty1981}
E. Cs\'aki and S. G. Mohanty,
Excursion and meander in random walk,
\emph{Canad. J. Statist.} \textbf{9} (1981), 57--70.
\href{https://doi.org/10.2307/3315296}{doi:10.2307/3315296}.

\bibitem{DurrettIglehartMiller1977}
R. T. Durrett, D. L. Iglehart, and D. R. Miller,
Weak convergence to Brownian meander and Brownian excursion,
\emph{Ann. Probab.} \textbf{5} (1977), 117--129.
\href{https://doi.org/10.1214/aop/1176995895}{doi:10.1214/aop/1176995895}.

\bibitem{HelfgottRadziwill2021}
H. A. Helfgott and M. Radziwi\l\l,
Expansion, divisibility and parity,
arXiv:2103.06853, 2021.

\bibitem{Kaigh1976}
W. D. Kaigh,
An invariance principle for random walk conditioned by a late return to zero,
\emph{Ann. Probab.} \textbf{4} (1976), 115--121.
\href{https://doi.org/10.1214/aop/1176996189}{doi:10.1214/aop/1176996189}.

\bibitem{KhorunzhiyMarckert2009}
O. Khorunzhiy and J.-F. Marckert,
Uniform bounds for exponential moment of maximum of a Dyck path,
\emph{Electron. Commun. Probab.} \textbf{14} (2009), 327--333.
\href{https://doi.org/10.1214/ECP.v14-1486}{doi:10.1214/ECP.v14-1486};
arXiv:0902.1229.

\bibitem{Kousha2012}
T. Kousha,
Asymptotic behavior and the moderate deviation principle for the maximum of a Dyck path,
\emph{Statist. Probab. Lett.} \textbf{82} (2012), 340--347.
\href{https://doi.org/10.1016/j.spl.2011.10.015}{doi:10.1016/j.spl.2011.10.015};
arXiv:1008.0606.

\bibitem{KoukoulopoulosSoundararajan2020}
D. Koukoulopoulos and K. Soundararajan,
The structure of multiplicative functions with small partial sums,
\emph{Discrete Anal.} (2020), Paper No. 6, 19 pp.
\href{https://doi.org/10.19086/da.11963}{doi:10.19086/da.11963};
arXiv:1909.13105.

\bibitem{Littlewood1912}
J. E. Littlewood,
Quelques cons\'equences de l'hypoth\`ese que la fonction $\zeta(s)$ de Riemann n'a pas de z\'eros dans le demi-plan
$\Re s>1/2$,
\emph{C. R. Acad. Sci. Paris} \textbf{154} (1912), 263--266.

\bibitem{MatomakiRadziwill2016}
K. Matom\"aki and M. Radziwi\l\l,
Multiplicative functions in short intervals,
\emph{Ann. of Math.} (2) \textbf{183} (2016), no.~3, 1015--1056.
\href{https://doi.org/10.4007/annals.2016.183.3.6}{doi:10.4007/annals.2016.183.3.6};
arXiv:1501.04585.

\bibitem{MatomakiRadziwillTao2015}
K. Matom\"aki, M. Radziwi\l\l, and T. Tao,
An averaged form of Chowla's conjecture,
\emph{Algebra Number Theory} \textbf{9} (2015), no.~9, 2167--2196.
\href{https://doi.org/10.2140/ant.2015.9.2167}{doi:10.2140/ant.2015.9.2167};
arXiv:1503.05121.

\bibitem{MatomakiRadziwillTao2016}
K. Matom\"aki, M. Radziwi\l\l, and T. Tao,
Sign patterns of the Liouville and M\"obius functions,
\emph{Forum Math. Sigma} \textbf{4} (2016), e14.
\href{https://doi.org/10.1017/fms.2016.6}{doi:10.1017/fms.2016.6};
arXiv:1509.01545.

\bibitem{MatomakiTeravainen2023}
K. Matom\"aki and J. Ter\"av\"ainen,
On the M\"obius function in all short intervals,
\emph{J. Eur. Math. Soc.} \textbf{25} (2023), no.~4, 1207--1225.
\href{https://doi.org/10.4171/JEMS/1205}{doi:10.4171/JEMS/1205};
arXiv:1911.09076.

\bibitem{MatomakiShaoTaoTeravainen2023}
K. Matom\"aki, X. Shao, T. Tao, and J. Ter\"av\"ainen,
Higher uniformity of arithmetic functions in short intervals I. All intervals,
\emph{Forum Math. Pi} \textbf{11} (2023), e29.
\href{https://doi.org/10.1017/fmp.2023.28}{doi:10.1017/fmp.2023.28};
arXiv:2204.03754.

\bibitem{Mirsky1949}
L. Mirsky,
On the frequency of pairs of square-free numbers with a given difference,
\emph{Bull. Amer. Math. Soc.} \textbf{55} (1949), 936--939.
\href{https://doi.org/10.1090/S0002-9904-1949-09313-8}{doi:10.1090/S0002-9904-1949-09313-8}.

\bibitem{Ng2004}
N. Ng,
The distribution of the summatory function of the M\"obius function,
\emph{Proc. London Math. Soc.} (3) \textbf{89} (2004), 361--389.
\href{https://doi.org/10.1112/S0024611504014741}{doi:10.1112/S0024611504014741};
arXiv:math/0310381.

\bibitem{Nunes2015}
R. M. Nunes,
Squarefree numbers in arithmetic progressions,
\emph{J. Number Theory} \textbf{153} (2015), 1--36.
\href{https://doi.org/10.1016/j.jnt.2014.12.025}{doi:10.1016/j.jnt.2014.12.025}.

\bibitem{OdlyzkoTeRiele1985}
A. M. Odlyzko and H. J. J. te Riele,
Disproof of the Mertens conjecture,
\emph{J. Reine Angew. Math.} \textbf{357} (1985), 138--160.
\href{https://doi.org/10.1515/crll.1985.357.138}{doi:10.1515/crll.1985.357.138}.

\bibitem{Pilatte2024}
C. Pilatte,
Improved bounds for the two-point logarithmic Chowla conjecture,
\emph{J. Amer. Math. Soc.}, in press.
\href{https://doi.org/10.1090/jams/1084}{doi:10.1090/jams/1084};
arXiv:2310.19357.

\bibitem{Tao2016}
T. Tao,
The logarithmically averaged Chowla and Elliott conjectures for two-point correlations,
\emph{Forum Math. Pi} \textbf{4} (2016), e8, 36 pp.
\href{https://doi.org/10.1017/fmp.2016.6}{doi:10.1017/fmp.2016.6};
arXiv:1509.05422.

\bibitem{TaoTeravainen2018}
T. Tao and J. Ter\"av\"ainen,
Odd order cases of the logarithmically averaged Chowla conjecture,
\emph{J. Th\'eor. Nombres Bordeaux} \textbf{30} (2018), no.~3, 997--1015.
\href{https://doi.org/10.5802/jtnb.1062}{doi:10.5802/jtnb.1062};
arXiv:1710.02112.

\bibitem{Titchmarsh1986}
E. C. Titchmarsh,
\emph{The Theory of the Riemann Zeta-Function},
2nd ed., revised by D. R. Heath-Brown, Clarendon Press, Oxford, 1986.

\end{thebibliography}

\end{document}
